\section{Behavioural Cloning - Double Pendulum Case}
\subsection{Behavioural Cloning model and policy training}
We trained a neural network with a horizon of $7$ actions ahead of the current one. Therefore the policy the agent attempts to learn is:

\begin{equation}
    \pi = (a_{t:t+7}|s_t)
\end{equation}

To train our model we first preproccess the expert data and instead of the state $s_t$ we input the feature transformation $\phi(s_t)$. Additionally, we separated the expert data into training and validation data. The training data comprise $90\%$ of the expert data  while testing data comprise $10\%$ of the expert data. We have performed normalisation in both the actions $a_t$ and the velocities $\dot{q}_t$ of the pendulum. In order to train this model we did not utilize test data since we test it in the simulation environment.

The architecture of the NN is depicted in Figure \ref{fig:NN-BC}

\begin{figure}[h]
    \centering
    \includegraphics[width=0.9\textwidth]{neural network.png}
    \caption{Architecture of BC policy model for the double pendulum example.}
    \label{fig:NN-BC}
\end{figure}

As it shown above the network has an input layer of size $6$ and 4 hidden fully connected layers while the output of the output layer is $16$ which is then reshaped to represent a set of $8$ consecutive actions in the form of a $8 \times 2$ matrix. The activation function of each layer is the hyperbolic tangent ($tanh(x)$). Additionally, in the output of each hidden layer we have added a dropout layer with a dropout rate of $20\%$ which is only active during training. Regarding the training hyperparameters are depicted in Table \ref{tab:bc-hyperparameters}. In Figure \ref{fig:bc-training} the training and test losses are depicted with respect to the training epochs. The training finishes early so as to not overfit to the expert trajectories.

\begin{figure}[h]
    \centering
    \includegraphics[width=0.7\textwidth]{training_bc.png}
    \caption{Train and Test loss curves of the BC policy model.}
    \label{fig:bc-training}
\end{figure}

\begin{table}[h]
  \centering
  \caption{Behavioural Cloning policy model and training hyperparameters (defaults).}\label{tab:bc-hyperparameters}
  \begin{tabular}{@{}lll@{}}
    \toprule
    \textbf{Parameter} & \textbf{Symbol} & \textbf{Value} \\
    \midrule
    Feature dimension         & ---            & $6$ per state \\
    Action prediction horizon & $H$            & $8$ \\
    Optimizer / LR            & ---            & Adam, $10^{-3}$ (cosine) \\
    Batch size / epochs       & ---            & $512\,/\,400$ \\
    Loss Function  & --- & MSE\\
    \bottomrule
  \end{tabular}
\end{table}



